\section*{第 \ref{ch:g3:sorting-living-things} 章——生物分类}

\begin{solution}{exo:g3:sorting-living-things:1}
\omterm{def:g3:sorting-living-things:criterion}{标准}是\omterm{def:g3:sorting-living-things:criterion}{分类}时所依据的规则。好的\omterm{def:g3:sorting-living-things:criterion}{标准}是可以观察到的身体特征，
如“有\omterm{def:g1:animals-around-us:covering}{羽毛}”“有六条腿”。
不好的\omterm{def:g3:sorting-living-things:criterion}{标准}是\omterm{def:g1:living-or-not:living}{生物}住在哪里或正在做什么，
如“住在水边”“会飞”。
\end{solution}

\begin{solution}{exo:g3:sorting-living-things:2}
\omterm{prop:g3:sorting-living-things:groups}{昆虫}：有六条腿，身体分三部分。\omterm{prop:g3:sorting-living-things:groups}{鱼类}：有鳍和鳞片，生活在水中。
\omterm{prop:g3:sorting-living-things:groups}{鸟类}：有\omterm{def:g1:animals-around-us:covering}{羽毛}、翅膀和喙。
\omterm{prop:g3:sorting-living-things:groups}{哺乳动物}：有毛发，雌性用\omterm{def:g2:animals-grow-up:born}{乳汁}喂养幼体。
\end{solution}

\begin{solution}{exo:g3:sorting-living-things:3}
蜜蜂：\omterm{prop:g3:sorting-living-things:groups}{昆虫}。金鱼：\omterm{prop:g3:sorting-living-things:groups}{鱼类}。鸵鸟：\omterm{prop:g3:sorting-living-things:groups}{鸟类}。奶牛：\omterm{prop:g3:sorting-living-things:groups}{哺乳动物}。青蛙：
\omterm{prop:g3:sorting-living-things:groups}{两栖动物}。
\end{solution}

\begin{solution}{exo:g3:sorting-living-things:4}
它有毛皮，用\omterm{def:g2:animals-grow-up:born}{乳汁}哺育幼体，没有\omterm{def:g1:animals-around-us:covering}{羽毛}，也没有喙。
飞行是它的一种行为，而不是定义\omterm{prop:g3:sorting-living-things:groups}{鸟类}的特征——
这些特征表明它是\omterm{prop:g3:sorting-living-things:groups}{哺乳动物}。
\end{solution}

\begin{solution}{exo:g3:sorting-living-things:5}
它没有鳞片，也没有\omterm{prop:g3:sorting-living-things:groups}{鱼类}那种鳃和鳍的身体结构——
它用\omterm{def:g2:animals-grow-up:born}{乳汁}哺育幼崽。游泳只是它的习性，
它的特征表明它是\omterm{prop:g3:sorting-living-things:groups}{哺乳动物}。
\end{solution}

\begin{solution}{exo:g3:sorting-living-things:6}
不对。\omterm{prop:g3:sorting-living-things:groups}{昆虫}的\omterm{def:g3:sorting-living-things:criterion}{标准}是六条腿（以及三部分身体）；蜘蛛有八条腿，
不符合这一\omterm{def:g3:sorting-living-things:criterion}{标准}。蜘蛛属于另一类群，
不论乍看起来像什么。
\end{solution}

\begin{solution}{exo:g3:sorting-living-things:7}
雏菊：有花和\omterm{def:g2:from-seed-to-plant:seed}{种子}，茎柔软、呈绿色。苔藓：从不开花，也没有真正的茎。
松树：是一棵树（有木质树干），叶呈针状，不开花，但长有带\omterm{def:g2:from-seed-to-plant:seed}{种子}的球果。
玫瑰丛：有花和\omterm{def:g2:from-seed-to-plant:seed}{种子}，
茎木质化。
\end{solution}

\begin{solution}{exo:g3:sorting-living-things:8}
开放性答案。大多数街道上\omterm{prop:g3:sorting-living-things:groups}{昆虫}数量遥遥领先——蚂蚁、
苍蝇、蜜蜂、甲虫、蝴蝶——\omterm{prop:g3:sorting-living-things:groups}{鸟类}排第二。
\end{solution}

\begin{solution}{exo:g3:sorting-living-things:9}
朋友的\omterm{def:g3:sorting-living-things:criterion}{分类}是“宠物与野生”。这种\omterm{def:g3:sorting-living-things:criterion}{分类}按\omterm{def:g1:animals-around-us:animal}{动物}与人类共处的方式划分，而这可能会改变——
比如一只流浪猫，无论在室内还是室外，都是同一种\omterm{def:g1:animals-around-us:animal}{动物}。
“\omterm{prop:g3:sorting-living-things:groups}{哺乳动物}与\omterm{prop:g3:sorting-living-things:groups}{鸟类}”则按一生都不变的身体特征划分，
这才是生物学家所用的那类\omterm{def:g3:sorting-living-things:criterion}{标准}。
\end{solution}

\begin{solution}{exo:g3:sorting-living-things:10}
第一步：覆盖——\omterm{def:g1:animals-around-us:covering}{羽毛}（紧密排列）、喙、两条腿、
翅膀（作鳍状）。第二步：幼体从卵中孵出。第三步：
特征与\omterm{prop:g3:sorting-living-things:groups}{鸟类}相符。第四步：“像鱼一样游泳、不会飞”的印象
与之不符，结果被否定。企鹅是\omterm{prop:g3:sorting-living-things:groups}{鸟类}。
\end{solution}
